\EMsection{The Two-Dimensional Wave Equation}{The Two-Dimensional Wave Equation}

Modelling a vibrating membrane (deriving the equation governing \EMim{u(x,y,t)}).

\textbf{Physical assumptions:}

\begin{itemize}
\tightlist
\item
  The mass is distributed uniformly over the membrane (a homogeneous membrane), and the membrane is perfectly flexible.
\item
  The membrane is stretched so tightly that the force of gravity on it is negligible.
\item
  The deflection of the membrane from equilibrium is very small.
\end{itemize}

\EMfigure{m13-membrane}{An element \EMim{\Delta x\times\Delta y} of the vibrating membrane and
the tensile forces \EMim{T\Delta x} and \EMim{T\Delta y} acting on it}

By a modelling process analogous to the one-dimensional case, the equation of the vibrating two-dimensional membrane is obtained:
\EMdisp{\frac{\partial^2u}{\partial t^2}\EMbr =c^2\left(\frac{\partial^2u}{\partial x^2}+\frac{\partial^2u}{\partial y^2}\right)\EMbr =c^2\nabla^2u}
For a unique solution we need boundary conditions and initial conditions:

\begin{itemize}
\tightlist
\item
  Boundary conditions: \EMim{u(x,y,t)=0} on the boundary of the membrane
\item
  Initial conditions: \EMim{u(x,y,0)=f(x,y)} and \EMim{u_t(x,y,0)=g(x,y)}, where \EMim{u_t\triangleq\dfrac{\partial u}{\partial t}}
\end{itemize}

\EMkeepfig{m13-rectangle}{3}

\EMsection{The Two-Dimensional Wave Equation with a Rectangular Boundary}{The Two-Dimensional Wave Equation with a Rectangular Boundary}

\EMfigure{m13-rectangle}{The rectangular region \EMim{0\le x\le a}, \EMim{0\le y\le b}}

\EMdisp{\frac{\partial^2u}{\partial t^2}\EMbr =c^2\left(\frac{\partial^2u}{\partial x^2}+\frac{\partial^2u}{\partial y^2}\right)}

\begin{itemize}
\tightlist
\item
  Boundary conditions: \EMim{u(0,y,t)=u(a,y,t)=0} and \EMim{u(x,0,t)=u(x,b,t)=0}
\item
  Initial conditions: \EMim{u(x,y,0)=f(x,y)} and \EMim{u_t(x,y,0)=g(x,y)}
\end{itemize}

\EMsubsection{Solving the Equation}{Solving the Equation}

Here we solve the equation by \textbf{separation of variables}. Assume \EMim{u(x,y,t)=X(x)\times Y(y)\times T(t)}.
\EMdisp{\frac{\partial^2u}{\partial x^2}\EMbr =X''YT,\qquad\EMbr \frac{\partial^2u}{\partial y^2}\EMbr =XY''T,\qquad\EMbr \frac{\partial^2u}{\partial t^2}\EMbr =XYT''\quad\EMbr \EMbr \Longrightarrow\quad\EMbr  XYT''\EMbr =c^2\big(X''YT+XY''T\big)}
\EMdisp{\frac{T''}{c^2T}\EMbr =\frac{X''}X+\frac{Y''}Y\EMbr =-\lambda^2\ \EMbr \Rightarrow\ \frac{Y''}Y\EMbr =-\lambda^2-\frac{X''}X\EMbr =-\rho^2\ \EMbr \Rightarrow\ \frac{X''}X\EMbr =-\lambda^2+\rho^2\EMbr =-\mu^2}
\EMdisp{\begin{cases}X''+\mu^2X=0\\ Y''+\rho^2Y=0\\ T''+c^2\lambda^2T=0\end{cases}\qquad\EMbr \lambda^2\EMbr =\mu^2+\rho^2}
\textbf{The solution in the \EMim{x} direction:}
\EMdisp{X(x)\EMbr =A_1\cos\mu x+A_2\sin\mu x,\quad\EMbr  X(0)\EMbr =0\EMbr \Rightarrow A_1\EMbr =0,\quad\EMbr  X(a)\EMbr =0\EMbr \Rightarrow A_2\sin\mu a\EMbr =0\EMbr \Rightarrow\mu a\EMbr =n\pi\EMbr \Rightarrow\mu_n\EMbr =\frac{n\pi}a}
\EMdisp{X_n(x)\EMbr =A_n\sin\mu_nx,\qquad\EMbr  \mu_n\EMbr =\frac{n\pi}a}
\textbf{The solution in the \EMim{y} direction:}
\EMdisp{Y(y)\EMbr =B_1\cos\rho y+B_2\sin\rho y,\quad\EMbr  Y(0)\EMbr =0\EMbr \Rightarrow B_1\EMbr =0,\quad\EMbr  Y(b)\EMbr =0\EMbr \Rightarrow B_2\sin\rho b\EMbr =0\EMbr \Rightarrow\rho b\EMbr =m\pi\EMbr \Rightarrow\rho_m\EMbr =\frac{m\pi}b}
\EMdisp{Y_m(y)\EMbr =B_m\sin\rho_my,\qquad\EMbr  \rho_m\EMbr =\frac{m\pi}b}
\textbf{The solution in time:}
\EMdisp{\lambda_{nm}^2\EMbr =\mu_n^2+\rho_m^2\EMbr =\left(\frac{n\pi}a\right)^2+\left(\frac{m\pi}b\right)^2,\qquad\EMbr  T_{nm}(t)\EMbr =C_1\cos(c\lambda_{nm}t)+C_2\sin(c\lambda_{nm}t)}
\EMdisp{u_{nm}(x,y,t)\EMbr =\big[A_{nm}\cos(c\lambda_{nm}t)+B_{nm}\sin(c\lambda_{nm}t)\big]\sin\!\left(\frac{n\pi}ax\right)\sin\!\left(\frac{m\pi}by\right),\qquad\EMbr  n,m\EMbr =1,2,\dots}
\EMdisp{u(x,y,t)\EMbr =\sum_{n=1}^{\infty}\sum_{m=1}^{\infty}\big[A_{nm}\cos(c\lambda_{nm}t)+B_{nm}\sin(c\lambda_{nm}t)\big]\sin\!\left(\frac{n\pi}ax\right)\sin\!\left(\frac{m\pi}by\right)}

\textbf{Applying the initial condition \EMim{u(x,y,0)=f(x,y)}:}
\EMdisp{f(x,y)\EMbr =\sum_{n=1}^{\infty}\sum_{m=1}^{\infty}A_{nm}\sin\!\left(\frac{n\pi}ax\right)\sin\!\left(\frac{m\pi}by\right)\EMbr =\sum_{n=1}^{\infty}K_n(y)\sin\frac{n\pi}ax,\qquad\EMbr  K_n(y)\triangleq\sum_{m=1}^{\infty}A_{nm}\sin\frac{m\pi}by}
\EMdisp{K_n(y)\EMbr =\frac2a\int_0^af(x,y)\sin\frac{n\pi}ax\,dx,\qquad\EMbr  A_{nm}\EMbr =\frac2b\int_0^bK_n(y)\sin\frac{m\pi}by\,dy}
\EMdisp{A_{nm}\EMbr =\frac4{ab}\int_0^b\left\{\int_0^af(x,y)\sin\frac{n\pi}ax\,dx\right\}\sin\frac{m\pi}by\,dy}

\begin{EMexercise}{}

Prove that
\EMdisp{B_{nm}\EMbr =\frac4{abc\lambda_{nm}}\int_0^b\left\{\int_0^ag(x,y)\sin\frac{n\pi}ax\,dx\right\}\sin\frac{m\pi}by\,dy}

\end{EMexercise}

\begin{EMimportant}{Solution of the two-dimensional wave equation on a rectangle}

\EMdisp{u(x,y,t)\EMbr =\sum_{n=1}^{\infty}\sum_{m=1}^{\infty}\big[A_{nm}\cos(c\lambda_{nm}t)+B_{nm}\sin(c\lambda_{nm}t)\big]\sin\!\left(\frac{n\pi}ax\right)\sin\!\left(\frac{m\pi}by\right),\qquad\EMbr \lambda_{nm}^2\EMbr =\left(\frac{n\pi}a\right)^2+\left(\frac{m\pi}b\right)^2}
\EMdisp{A_{nm}\EMbr =\frac4{ab}\int_0^b\left\{\int_0^af(x,y)\sin\frac{n\pi}ax\,dx\right\}\sin\frac{m\pi}by\,dy,\qquad\EMbr  B_{nm}\EMbr =\frac4{abc\lambda_{nm}}\int_0^b\left\{\int_0^ag(x,y)\sin\frac{n\pi}ax\,dx\right\}\sin\frac{m\pi}by\,dy}

\end{EMimportant}

\begin{EMexample}{}

Find the solution of the two-dimensional motion equation above for \EMim{g(x,y)=0} and \EMim{f(x,y)=xy}.

\EMdisp{g(x,y)\EMbr =0\ \EMbr \Rightarrow\ B_{nm}\EMbr =0}
\EMdisp{A_{nm}\EMbr =\frac4{ab}\int_0^b\left\{\int_0^axy\sin\frac{n\pi}ax\,dx\right\}\sin\frac{m\pi}by\,dy\EMbr =\frac4{ab}\left\{\left[\int_0^by\sin\frac{m\pi}by\,dy\right]\times\left[\int_0^ax\sin\frac{n\pi}ax\,dx\right]\right\}}
\EMdisp{=\frac4{ab}\left[y\left(\frac{-b}{m\pi}\cos\frac{m\pi}by\right)-(1)\left(\frac{-b^2}{m^2\pi^2}\sin\frac{m\pi}by\right)\right]_0^b\times\left[x\left(\frac{-a}{n\pi}\cos\frac{n\pi}ax\right)-(1)\left(\frac{-a^2}{n^2\pi^2}\sin\frac{n\pi}ax\right)\right]_0^a}
\EMdisp{=\frac4{ab}\cdot\frac{-b^2}{m\pi}(-1)^m\cdot\frac{-a^2}{n\pi}(-1)^n\EMbr =\frac{4ab}{nm\pi^2}(-1)^{n+m}}
\EMdisp{u(x,y,t)\EMbr =\sum_{n=1}^{\infty}\sum_{m=1}^{\infty}\frac{4ab}{nm\pi^2}(-1)^{n+m}\cos(c\lambda_{nm}t)\sin\!\left(\frac{n\pi}ax\right)\sin\!\left(\frac{m\pi}by\right)}

\end{EMexample}

\EMsection{The Two-Dimensional Heat Equation with a Rectangular Boundary}{The Two-Dimensional Heat Equation with a Rectangular Boundary}

\EMdisp{\frac{\partial u}{\partial t}\EMbr =c^2\left(\frac{\partial^2u}{\partial x^2}+\frac{\partial^2u}{\partial y^2}\right)\EMbr =c^2\nabla^2u}

\begin{itemize}
\tightlist
\item
  Boundary condition: \EMim{u(x,y,t)=0} on the boundary of the plate
\item
  Initial condition: \EMim{u(x,y,0)=f(x,y)}
\end{itemize}

\EMsubsection{Solving the Equation}{Solving the Equation}

Assume \EMim{u(x,y,t)=X(x)\times Y(y)\times T(t)}.
\EMdisp{\frac{\partial^2u}{\partial x^2}\EMbr =X''YT,\qquad\EMbr \frac{\partial^2u}{\partial y^2}\EMbr =XY''T,\qquad\EMbr \frac{\partial u}{\partial t}\EMbr =XYT'\quad\EMbr \EMbr \Longrightarrow\quad\EMbr  XYT'\EMbr =c^2\big(X''YT+XY''T\big)}
\EMdisp{\frac{T'}{c^2T}\EMbr =\frac{X''}X+\frac{Y''}Y\EMbr =-\lambda^2,\qquad\EMbr  \frac{Y''}Y\EMbr =-\lambda^2-\frac{X''}X\EMbr =-\rho^2,\qquad\EMbr \frac{X''}X\EMbr =-\lambda^2+\rho^2\EMbr =-\mu^2}
\EMdisp{\begin{cases}X''+\mu^2X=0\\ Y''+\rho^2Y=0\\ T'+c^2\lambda^2T=0\end{cases}\qquad\EMbr \lambda^2\EMbr =\mu^2+\rho^2}
As in the wave case (\EMim{X_n(x)=A_n\sin\mu_nx} with \EMim{\mu_n=n\pi/a} and \EMim{Y_m(y)=B_m\sin\rho_my} with \EMim{\rho_m=m\pi/b}), and with
\EMdisp{T_{nm}(t)\EMbr =C_1e^{-c^2\lambda_{nm}^2t},\qquad\EMbr \lambda_{nm}^2\EMbr =\mu_n^2+\rho_m^2\EMbr =\left(\frac{n\pi}a\right)^2+\left(\frac{m\pi}b\right)^2}
\EMdisp{u_{nm}(x,y,t)\EMbr =A_{nm}\sin\!\left(\frac{n\pi}ax\right)\sin\!\left(\frac{m\pi}by\right)e^{-c^2\lambda_{nm}^2t},\qquad\EMbr  n\EMbr =1,2,\dots,\ m\EMbr =1,2,\dots}
\EMdisp{u(x,y,t)\EMbr =\sum_{n=1}^{\infty}\sum_{m=1}^{\infty}A_{nm}\sin\!\left(\frac{n\pi}ax\right)\sin\!\left(\frac{m\pi}by\right)e^{-c^2\lambda_{nm}^2t}}
\textbf{Applying the initial condition:} exactly as in the wave case,
\EMdisp{A_{nm}\EMbr =\frac4{ab}\int_0^b\left\{\int_0^af(x,y)\sin\frac{n\pi}ax\,dx\right\}\sin\frac{m\pi}by\,dy}

\EMkeepfig{m13-polar}{5}

\EMsection{The Laplacian in Polar Coordinates}{The Laplacian in Polar Coordinates}

The aim is to write \EMim{\nabla^2u=u_{xx}+u_{yy}} in polar coordinates in terms of \EMim{r,\theta}.

\EMfigure{m13-polar}{Polar coordinates: \EMim{x=r\cos\theta}, \EMim{y=r\sin\theta}}

\EMdisp{r\EMbr =\sqrt{x^2+y^2},\quad\EMbr  x\EMbr =r\cos\theta,\qquad\EMbr \theta\EMbr =\tan^{-1}\!\left(\frac yx\right),\quad\EMbr  y\EMbr =r\sin\theta}
\EMdisp{u_x\EMbr =u_rr_x+u_\theta\theta_x\ \EMbr \Rightarrow\ u_{xx}\EMbr =\big(u_rr_x+u_\theta\theta_x\big)_x\EMbr =u_{rr}r_x^2+u_{\theta\theta}\theta_x^2+2r_x\theta_xu_{r\theta}+r_{xx}u_r+\theta_{xx}u_\theta}
\EMdisp{u_y\EMbr =u_rr_y+u_\theta\theta_y\ \EMbr \Rightarrow\ u_{yy}\EMbr =u_{rr}r_y^2+u_{\theta\theta}\theta_y^2+2r_y\theta_yu_{r\theta}+r_{yy}u_r+\theta_{yy}u_\theta}
\EMdisp{r_x\EMbr =\frac xr,\quad\EMbr  r_{xx}\EMbr =\frac{y^2}{r^3},\quad\EMbr \theta_x\EMbr =\frac{-y}{r^2},\quad\EMbr \theta_{xx}\EMbr =\frac{2xy}{r^4},\qquad\EMbr  r_y\EMbr =\frac yr,\quad\EMbr  r_{yy}\EMbr =\frac{x^2}{r^3},\quad\EMbr \theta_y\EMbr =\frac x{r^2},\quad\EMbr \theta_{yy}\EMbr =\frac{-2xy}{r^4}}
\EMdisp{u_{xx}\EMbr =\frac{x^2}{r^2}u_{rr}+\frac{y^2}{r^4}u_{\theta\theta}+\frac{-2xy}{r^3}u_{r\theta}+\frac{y^2}{r^3}u_r+\frac{2xy}{r^4}u_\theta}
\EMdisp{u_{yy}\EMbr =\frac{y^2}{r^2}u_{rr}+\frac{x^2}{r^4}u_{\theta\theta}+\frac{2xy}{r^3}u_{r\theta}+\frac{x^2}{r^3}u_r+\frac{-2xy}{r^4}u_\theta}
Adding these two relations:

\begin{EMimportant}{The Laplacian in polar coordinates}

\EMdisp{\nabla^2u\EMbr =u_{rr}+\frac1{r^2}u_{\theta\theta}+\frac1ru_r}

\end{EMimportant}

\EMkeepfig{m13-disc}{3}

\EMsection{The Two-Dimensional Wave Equation with a Circular Boundary}{The Two-Dimensional Wave Equation with a Circular Boundary}

\EMfigure{m13-disc}{A circular membrane of radius \EMim{R}}

\EMdisp{\frac{\partial^2u}{\partial t^2}\EMbr =c^2\nabla^2u,\qquad\EMbr  u(R,t)\EMbr =0\ \ \text{(boundary condition)},\qquad\EMbr  u(r,0)\EMbr =f(r),\ \ u_t(r,0)\EMbr =g(r)\ \ \text{(initial conditions)}}

\EMsubsection{Solving the Equation}{Solving the Equation}

Because of the circular symmetry of this problem it is better to solve it in polar coordinates. On a circular membrane the vibration is not a function of \EMim{\theta}, so
\EMdisp{\begin{cases}\dfrac{\partial^2u}{\partial t^2}=c^2\left(u_{rr}+\dfrac1ru_r\right)\\\noalign{\vskip 2mm} u(R,t)=0,\quad u(r,0)=f(r),\quad u_t(r,0)=g(r)\end{cases}}
Now, assuming separable variables \EMim{u(r,t)=P(r)\times T(t)}:
\EMdisp{PT''\EMbr =c^2\left(P''T+\frac1rP'T\right)\ \EMbr \Rightarrow\ \frac{T''}{c^2T}\EMbr =\frac{P''}P+\frac1r\frac{P'}P\EMbr =-\lambda^2\ \EMbr \Rightarrow\ \begin{cases}P''+\dfrac1rP'+\lambda^2P=0\\\noalign{\vskip 1mm} T''+c^2\lambda^2T=0\end{cases}}

\EMkeepfig{m13-bessel}{5}

\EMsubsection{Bessel's Differential Equation of Order \EMim{n}}{Bessel’s Differential Equation of Order n}

\EMdisp{y''+\frac1xy'+\left(1-\frac{n^2}{x^2}\right)y\EMbr =0\quad\EMbr \EMbr \Longrightarrow\quad\EMbr  y(x)\EMbr =AJ_n(x)+BY_n(x)}
where \EMim{J_n} are the Bessel functions of the first kind of order \EMim{n} and \EMim{Y_n} the Bessel functions of the second kind of order \EMim{n}.

\EMfigure{m13-bessel}{The Bessel functions of the first kind \EMim{J_0,\dots,J_4} and of the
second kind \EMim{Y_0,\dots,Y_4} for \EMim{0<x\le10}}

With the substitution \EMim{s=\lambda r}:
\EMdisp{P'\EMbr =\frac{dP}{dr}\EMbr =\frac{dP}{ds}\frac{ds}{dr}\EMbr =\lambda\frac{dP}{ds},\qquad\EMbr  P''\EMbr =\lambda^2\frac{d^2P}{ds^2},\qquad\EMbr \frac1r\lambda\EMbr =\frac{\lambda^2}s}
\EMdisp{\lambda^2\frac{d^2P}{ds^2}+\frac1r\lambda\frac{dP}{ds}+\lambda^2P\EMbr =0\ \EMbr \Rightarrow\ \frac{d^2P}{ds^2}+\frac1s\frac{dP}{ds}+P\EMbr =0}
This is Bessel's equation of order zero (\EMim{n=0}), so
\EMdisp{P(r)\EMbr =AJ_0(\lambda r)+BY_0(\lambda r)}
\textbf{Applying the boundary conditions:}

\begin{itemize}
\tightlist
\item
  Since \EMim{P(0)\neq\infty} and \EMim{Y_0} is unbounded at the origin, the coefficient of \EMim{Y_0} is zero (\EMim{D=0}) and \EMim{P(r)=CJ_0(\lambda r)}.
\item
  \EMim{P(R)=0} gives \EMim{J_0(\lambda R)=0}, i.e.~\EMim{\lambda R=\alpha_m}, where \EMim{\alpha_m} are the roots of \EMim{J_0}:
  \EMdisp{\lambda\EMbr =\lambda_m\EMbr =\frac{\alpha_m}R,\quad\EMbr  m\EMbr =1,2,\dots\quad\EMbr \EMbr \Longrightarrow\quad\EMbr  P_m(r)\EMbr =C_mJ_0\!\left(\frac{\alpha_m}Rr\right)}
\end{itemize}

\EMfigure{m13-j0-zeros}{The function \EMim{J_0(x)} and its roots
\EMim{\alpha_1,\alpha_2,\alpha_3}}

\EMdisp{T''+c^2\lambda_m^2T\EMbr =0\ \EMbr \Rightarrow\ T_m(t)\EMbr =A_m\cos(c\lambda_mt)+B_m\sin(c\lambda_mt)}

\begin{EMimportant}{Solution of the two-dimensional wave equation on a disc}

\EMdisp{u(r,t)\EMbr =\sum_{m=1}^{\infty}\big(A_m\cos(c\lambda_mt)+B_m\sin(c\lambda_mt)\big)J_0\!\left(\frac{\alpha_m}Rr\right),\qquad\EMbr \lambda_m\EMbr =\frac{\alpha_m}R}

\end{EMimportant}

\begin{EMremark}{Research task (up to 1 bonus point)}

Prepare a complete report on Bessel's differential equation of order \EMim{n}, the Bessel functions of the first and second kind and their properties, as well as the complete solution of the two-dimensional wave equation with a circular boundary (the report must be typed, in Word format).

\end{EMremark}
